\documentclass[final]{beamer}

\usepackage[size=a0,orientation=landscape,scale=1.00]{beamerposter}
\usepackage[T1]{fontenc}
\usepackage{lmodern}
\usepackage{graphicx}
\usepackage{booktabs}
\usepackage{tabularx}
\usepackage{array}
\usepackage{ragged2e}
\usepackage{tikz}
\usetikzlibrary{arrows.meta,positioning,calc}

\setbeamertemplate{navigation symbols}{}
\setbeamertemplate{blocks}[rounded][shadow=false]
\setbeamersize{text margin left=1.18cm,text margin right=1.18cm}

\definecolor{Navy}{HTML}{173F5F}
\definecolor{DeepNavy}{HTML}{12344D}
\definecolor{Teal}{HTML}{238A8D}
\definecolor{Green}{HTML}{2E8B57}
\definecolor{Orange}{HTML}{C76B20}
\definecolor{BluePanel}{HTML}{F1F6F9}
\definecolor{GreenPanel}{HTML}{EFF8F3}
\definecolor{OrangePanel}{HTML}{FFF4EA}
\definecolor{GrayPanel}{HTML}{F5F6F7}
\definecolor{LineGray}{HTML}{CDD6DC}
\definecolor{BodyText}{HTML}{18242D}

\setbeamercolor{background canvas}{bg=white}
\setbeamercolor{normal text}{fg=BodyText,bg=white}
\setbeamercolor{block title}{fg=white,bg=Navy}
\setbeamercolor{block body}{fg=BodyText,bg=BluePanel}
\setbeamercolor{block title alerted}{fg=white,bg=Teal}
\setbeamercolor{block body alerted}{fg=BodyText,bg=GreenPanel}
\setbeamercolor{block title example}{fg=white,bg=Green}
\setbeamercolor{block body example}{fg=BodyText,bg=GreenPanel}

\setbeamerfont{block title}{size=\Large,series=\bfseries}
\setbeamerfont{block body}{size=\normalsize}

\newcommand{\StudentID}{40310663}

\newcommand{\metric}[3]{%
  \begin{tikzpicture}[baseline=(m.center)]
    \node[fill=BluePanel,rounded corners=2pt,minimum width=0.180\textwidth,
      minimum height=2.05cm,inner sep=0pt,align=center] (m) {%
      {\fontsize{29}{33}\selectfont\bfseries\color{#1}#2}\\[-0.01cm]
      {\fontsize{14}{17}\selectfont #3}%
    };
  \end{tikzpicture}%
}
\newcommand{\imgbox}[2]{%
  \noindent\makebox[\linewidth][c]{%
    \fcolorbox{LineGray}{white}{\includegraphics[width=0.975\linewidth,#1,keepaspectratio]{#2}}%
  }\par
}
\newcommand{\capline}[1]{%
  \vspace{0.06cm}%
  \noindent\parbox{\linewidth}{\centering\footnotesize\color{DeepNavy}#1}\par
  \vspace{0.08cm}%
}
\newcommand{\req}[2]{\textbf{#1} & #2\\}
\newcommand{\vdivider}[1]{%
  \begin{column}{0.007\textwidth}
  \centering\color{LineGray!85}%
  \raisebox{-#1}[0pt][0pt]{\rule{0.05cm}{#1}}%
  \end{column}%
}

\begin{document}
\begin{frame}[t]

% ========================= HEADER =========================
\begin{beamercolorbox}[wd=\textwidth,sep=0.86cm,leftskip=1.25cm,rightskip=1.25cm]{block title}
\begin{minipage}[c]{0.67\textwidth}
{\fontsize{57}{63}\selectfont\bfseries Gamma Calculator}\\[0.10cm]
{\fontsize{30}{36}\selectfont Deliverable 3 -- Code Quality, UIDP, Accessibility \& Unit Testing}\\[0.08cm]
{\fontsize{22}{27}\selectfont F4: $\Gamma(x)$ -- From-scratch Python scientific calculator}
\end{minipage}\hfill
\begin{minipage}[c]{0.30\textwidth}
\raggedleft
{\fontsize{24}{29}\selectfont\bfseries Md. Ariful Hossain}\\[0.08cm]
{\fontsize{18}{22}\selectfont \textbf{Student ID:} \StudentID}\\
{\fontsize{18}{22}\selectfont SOEN 6011: Software Engineering Processes}\\
{\fontsize{18}{22}\selectfont Summer 2026}\\[0.18cm]
{\fontsize{14}{18}\selectfont\textbf{Repository:} arifulHossain1/\\[-0.01cm]
soen6011-Scientific-Calculator-gamma-d2}
\end{minipage}
\end{beamercolorbox}

\vspace{0.28cm}
\centering
\noindent\begin{tabular}{@{}c@{\hspace{0.55cm}}c@{\hspace{0.55cm}}c@{\hspace{0.55cm}}c@{\hspace{0.55cm}}c@{}}
\metric{Teal}{0}{Flake8 violations} &
\metric{Navy}{9.91/10}{Pylint score} &
\metric{Green}{22/22}{PyUnit tests passed} &
\metric{Orange}{v1.1.0}{D3 semantic version} &
\metric{Teal}{$1.332\times10^{-15}$}{max relative error}
\end{tabular}
\par
\vspace{0.34cm}

% ========================= TOP ROW =========================
\begin{columns}[T,totalwidth=0.970\textwidth]

% -------- TOP LEFT: P7 QUALITY --------
\begin{column}{0.306\textwidth}
\begin{block}{D3 Problem 7 (70 marks) -- Code quality evidence}
\small
\renewcommand{\arraystretch}{1.13}
\begin{tabularx}{\linewidth}{>{\bfseries\RaggedRight}p{0.23\linewidth} X}
\req{PEP 8}{Production and test code formatted for naming, imports, spacing, wrapping and docstrings.}
\req{Flake8}{Style conformance checked from CMD; final run reports \textbf{0 violations}.}
\req{Debugger}{Python \texttt{pdb}: breakpoint, stepping, variable inspection and return-value inspection.}
\req{Pylint}{Static analysis completed with a final score of \textbf{9.91/10}.}
\req{SemVer}{D2 baseline \textbf{v1.0.0}; D3 backward-compatible feature release \textbf{v1.1.0}.}
\end{tabularx}
\end{block}

\vspace{0.20cm}
\begin{alertblock}{Flake8 -- enlarged PEP 8 evidence}
\imgbox{height=5.6cm}{assets/flake8.png}
\capline{Command: \texttt{python -m flake8 src tests --count --statistics}; final result: \textbf{0}.}
\end{alertblock}

\vspace{0.20cm}
\begin{block}{Pylint -- enlarged static-analysis evidence}
\imgbox{height=7.2cm}{assets/pylint.png}
\vspace{0.08cm}
\scriptsize
\noindent\textbf{Remaining Pylint messages reviewed and retained:}\par
\vspace{0.03cm}
\begin{itemize}\setlength\itemsep{0.02cm}\setlength\parskip{0pt}\setlength\parsep{0pt}
\item \texttt{R0912/R0915}: explicit parser branches/statements are retained because the from-scratch parser is intentionally transparent and traceable.
\item \texttt{R1731}: the explicit running-maximum \texttt{if} in validation is retained for readability; this is a refactoring suggestion, not a correctness defect.
\end{itemize}
\vspace{-0.03cm}
\noindent\textbf{Decision:} retain correct, requirement-aligned code rather than change it only to obtain 10/10.\par
\end{block}

\vspace{0.20cm}
\begin{block}{pdb -- what was debugged and what it confirmed}
\imgbox{height=10.5cm}{assets/pdb.png}
\vspace{0.08cm}
\scriptsize
\begin{tabularx}{\linewidth}{>{\bfseries}p{0.25\linewidth} X}
Breakpoint & \texttt{main()} in \texttt{tools/debug\_gamma.py} (line 12 in the captured run).\\
Commands & \texttt{b main}, \texttt{c}, \texttt{n}, \texttt{p x\_value}, \texttt{s}, \texttt{r}, \texttt{p result}, \texttt{q}.\\
Inspected & \texttt{x\_value=-0.5}; returned result \texttt{-3.5449077018110295}.\\
Conclusion & Confirmed that $x=-0.5$ enters \texttt{gamma\_lanczos()}, takes the $x<0.5$ Euler-reflection branch and returns the expected value. No defect was found in this session.\\
\end{tabularx}
\end{block}
\end{column}

\vdivider{45.0cm}

% -------- TOP CENTER: UIDP --------
\begin{column}{0.330\textwidth}
\begin{exampleblock}{D3 Problem 7 -- UIDP radial mind map}
\centering
\includegraphics[width=0.99\linewidth,height=27.2cm,keepaspectratio]{uidp_mindmap.pdf}
\vspace{0.10cm}
\scriptsize
\begin{tabularx}{0.98\linewidth}{>{\bfseries\RaggedRight}p{0.27\linewidth} X}
Applied directly & User Profiling, Metaphor, Feature Exposure, State Visualization, Coherence and Safety.\\
Partial only & Humility remains partial because no representative-user study was conducted.\\
Mind-map style & Central topic + seven branches + compact keyword cues for each principle. [2]\\
\end{tabularx}
\end{exampleblock}

\vspace{0.20cm}
\begin{block}{UIDP interpretation}
\small
\begin{itemize}\setlength\itemsep{0.10cm}
\item The GUI emphasizes a \textbf{simple calculator workflow}: one input, visible actions, readable output and recovery from invalid input.
\item \textbf{Metaphor} is used lightly through a familiar form/calculator structure; \textbf{Humility} stays partial because representative-user evidence was not collected.
\item The radial format improves readability: each principle is a branch and each branch uses only short cues rather than paragraph-style text.
\end{itemize}
\end{block}
\end{column}

\vdivider{45.0cm}

% -------- TOP RIGHT: P8 --------
\begin{column}{0.306\textwidth}
\begin{exampleblock}{D3 Problem 8 (30 marks) -- PyUnit test design}
\small
\renewcommand{\arraystretch}{1.12}
\begin{tabularx}{\linewidth}{>{\bfseries\RaggedRight}p{0.28\linewidth} X}
\req{Input parsing}{Decimal, negative decimal, scientific notation, empty input, alphabetic input and multiple values.}
\req{Gamma values}{$\Gamma(1)$, $\Gamma(5)$, $\Gamma(0.5)$, $\Gamma(-0.5)$ and $\Gamma(-1.5)$.}
\req{Domain / range}{Exact poles, negative non-integers, near-pole detection, domain exceptions and range overflow.}
\req{User-facing output}{Direct Lanczos path, reflection path, result text and near-pole warning.}
\req{Framework}{Python \texttt{unittest} (PyUnit), executed with test discovery from CMD.}
\end{tabularx}
\vspace{0.08cm}
\begin{center}
\fcolorbox{Green}{white}{\parbox{0.88\linewidth}{\centering\bfseries\Large 22 tests executed -- 22 passed -- 0 failures}}
\end{center}
\end{exampleblock}

\vspace{0.20cm}
\begin{block}{Representative unit-test cases and oracles}
\scriptsize
\renewcommand{\arraystretch}{1.17}
\begin{tabularx}{\linewidth}{>{\bfseries}p{0.17\linewidth} p{0.22\linewidth} p{0.25\linewidth} X}
\toprule
Test & Input & Expected & Test oracle\\
\midrule
Positive integer & $5$ & $24$ & Mathematical identity $\Gamma(n)=(n-1)!$.\\
Reflection case & $-0.5$ & $-2\sqrt{\pi}\approx-3.544907701811032$ & DLMF Gamma identities / reflection behavior.\\
Domain error & $0$ & \texttt{GammaDomainError} & Gamma has poles at nonpositive integers.\\
\bottomrule
\end{tabularx}
\vspace{0.08cm}
\textbf{Distinction:} these are part of the \textbf{22 automated unit tests}. They are separate from the independent eight-case numerical validation below.
\end{block}

\vspace{0.20cm}
\begin{block}{PyUnit execution evidence}
\imgbox{height=10.7cm}{assets/unittest.png}
\capline{Verbose discovery ends with \textbf{Ran 22 tests} and \textbf{OK}.}
\end{block}

\vspace{0.20cm}
\begin{alertblock}{Accuracy evidence -- separate eight-case validation}
\scriptsize
\textbf{Independent oracle:} Python standard-library \texttt{math.gamma()} in \texttt{src/validate\_gamma.py}; it is not imported by the production calculator. [11]
\par
\textbf{Validation set:} $\{0.1,0.5,1,1.5,2.5,5,-0.5,-1.5\}$. The reported \textbf{$1.332\times10^{-15}$} is the \textbf{maximum relative error over these eight cases}, below the required $10^{-10}$ threshold.
\vspace{0.08cm}
\imgbox{height=8.4cm}{assets/validation.png}
\capline{Separate regression validation: all eight reference cases PASS.}
\end{alertblock}
\end{column}
\end{columns}

\vspace{0.32cm}
\noindent\makebox[\textwidth][c]{\color{LineGray!85}\rule{0.970\textwidth}{0.075cm}}
\vspace{0.30cm}

% ========================= BOTTOM ROW =========================
\begin{columns}[T,totalwidth=0.970\textwidth]

% -------- BOTTOM LEFT: SEMVER + REPO --------
\begin{column}{0.306\textwidth}
\begin{block}{Semantic Versioning -- why v1.0.0 $\rightarrow$ v1.1.0}
\small
\begin{center}
\begin{tikzpicture}[>=Latex,node distance=0.55cm]
\node[draw=Navy,very thick,rounded corners=5pt,fill=BluePanel,text width=5.6cm,minimum height=1.75cm,align=center] (v1) {\Large\textbf{v1.0.0}\\\normalsize D2 stable baseline};
\node[draw=Teal,very thick,rounded corners=5pt,fill=GreenPanel,text width=8.1cm,minimum height=2.45cm,align=center,right=0.75cm of v1] (features) {\bfseries Backward-compatible user-visible D3 additions\\\scriptsize keyboard shortcuts; focus return/selection; clearer status/help text; scrollable history; accessibility improvements};
\node[draw=Orange,very thick,rounded corners=5pt,fill=OrangePanel,text width=5.6cm,minimum height=1.75cm,align=center,right=0.75cm of features] (v11) {\Large\textbf{v1.1.0}\\\normalsize D3 release};
\draw[->,very thick,Navy] (v1) -- (features);
\draw[->,very thick,Navy] (features) -- (v11);
\end{tikzpicture}
\end{center}
\begin{tabularx}{\linewidth}{>{\bfseries\RaggedRight}p{0.27\linewidth} X}
Why MINOR? & The release adds backward-compatible, user-visible GUI/accessibility features. Internal testing/tooling alone is \emph{not} the sole justification.\\
Why not PATCH? & More than a bug fix: new keyboard/focus/history/status behavior is visible to users.\\
Why not MAJOR? & No incompatible change to the D2 Gamma algorithm or expected calculation workflow.\\
\end{tabularx}
\vspace{0.07cm}
\imgbox{height=4.8cm}{assets/semantic_version.png}
\capline{Release evidence: \textbf{v1.0.0} preserves the D2 baseline; \textbf{v1.1.0} identifies the backward-compatible D3 feature release.}
\end{block}

\vspace{0.18cm}
\begin{block}{GitHub repository \& final source structure}
\scriptsize
\textbf{Repository:} \texttt{github.com/arifulHossain1/}\\[-0.02cm]
\hspace*{2.1cm}\texttt{soen6011-Scientific-Calculator-gamma-d2}
\vspace{0.07cm}
\begin{tabularx}{\linewidth}{>{\ttfamily\bfseries}p{0.37\linewidth} X}
src/gamma\_core.py & Parsing, numerical helpers, exceptions and Gamma calculation.\\
src/gui.py & Tkinter interface, keyboard/focus behavior, status and history.\\
src/main.py & IDE-independent entry point.\\
src/validate\_gamma.py & Separate numerical validation using \texttt{math.gamma()} as oracle.\\
src/version.py & Semantic version identifier (1.1.0).\\
tests/test\_gamma\_core.py & 22 PyUnit tests.\\
tools/debug\_gamma.py & pdb debugging driver.\\
\end{tabularx}
\end{block}
\end{column}

\vdivider{25.6cm}

% -------- BOTTOM CENTER: ACCESSIBILITY --------
\begin{column}{0.330\textwidth}
\begin{alertblock}{Accessibility Features and Evidence}
\small
\textbf{Scope:} the GUI \emph{aims} to be accessible; no formal representative-user, screen-reader or assistive-technology study is claimed.
\vspace{0.08cm}
\renewcommand{\arraystretch}{1.13}
\begin{tabularx}{\linewidth}{>{\bfseries\RaggedRight}p{0.31\linewidth} X}
Keyboard access & Enter / Alt+C = Calculate; Esc / Alt+L = Clear.\\
Focus management & Focus starts and returns to input; text is selected after success/error for rapid correction.\\
Readability & Explicit label, larger history font, scrollbar, visible status region.\\
No color-only meaning & Results, warnings and errors are communicated in text.\\
Safety / recovery & Invalid input produces a corrective message without terminating the program.\\
Metaphor & Familiar form/calculator structure: input field, Calculate/Clear buttons, status and history.\\
\end{tabularx}
\vspace{0.10cm}
\begin{columns}[T,totalwidth=\linewidth]
\begin{column}{0.495\linewidth}
\centering
\includegraphics[width=\linewidth,height=8.6cm,keepaspectratio]{assets/gui_success.png}\par
\vspace{0.05cm}
\noindent\parbox{\linewidth}{\centering\scriptsize\textbf{Successful use:} $x=-0.5$ shows the result, reflection method and completion status.}\par
\end{column}
\begin{column}{0.495\linewidth}
\centering
\includegraphics[width=\linewidth,height=8.6cm,keepaspectratio]{assets/gui_error.png}\par
\vspace{0.05cm}
\noindent\parbox{\linewidth}{\centering\scriptsize\textbf{Recoverable error:} invalid text gives a corrective message and visible status.}\par
\end{column}
\end{columns}
\end{alertblock}

\vspace{0.18cm}
\begin{block}{D2 numerical core preserved}
\scriptsize
D3 changes improve software quality and interaction without replacing the selected mathematics. Production still uses the \textbf{Lanczos approximation} with \textbf{Euler reflection} and project-defined subordinate numerical helpers. The independent \texttt{math.gamma()} oracle remains isolated to validation. [9--11]
\end{block}
\end{column}

\vdivider{25.6cm}

% -------- BOTTOM RIGHT: GAI / TRACEABILITY / REFS --------
\begin{column}{0.306\textwidth}
\begin{alertblock}{GAI / CASTROFF -- Problems 7 and 8}
\footnotesize
\textbf{Framework source:} CASTROFF = Constraints, Audience, Structure, Tone, Role, Output Format, Focus, Function. [12] \textbf{Tool:} ChatGPT. Tool screenshots and executed tests are the evidence; GAI was used for planning/review, not to fabricate results.
\vspace{0.08cm}

\textbf{Problem 7 -- quality, UIDP and accessibility}\\
\textbf{Compact prompt example:} ``Review my D2 Gamma Calculator for D3. Preserve the from-scratch numerical behavior; guide PEP 8/Flake8/Pylint/pdb/SemVer work, map the course UIDP principles, and improve accessibility with evidence suitable for the poster.''\\
\textbf{CASTROFF:} C = preserve D2 + meet P7; A = SOEN 6011 evaluator; S = stepwise code/tool/evidence workflow; T = academic/verifiable; R = Python quality + UI reviewer; O = revised code/commands/evidence; F = P7 quality/UIDP/accessibility; Function = defensible D3 release.\\
\textbf{Output summary:} code-quality workflow, accessibility changes, UIDP mapping and poster evidence plan.\\
\textbf{Decision:} \textcolor{Green}{\bfseries Accepted} PEP 8/tool workflow + keyboard/focus/readability changes; \textcolor{Teal}{\bfseries Revised} UIDP mapping and SemVer rationale after requirements review; \textcolor{Orange}{\bfseries Rejected} treating tool/test/reviewer evidence as representative-user input or changing correct code only to chase Pylint 10/10.

\vspace{0.10cm}
\textbf{Problem 8 -- unit testing}\\
\textbf{Compact prompt example:} ``Create PyUnit tests for the actual Gamma core covering parsing, direct Lanczos, Euler reflection, poles/near-poles, range errors, exceptions and user-facing output; keep independent numerical validation separate.''\\
\textbf{CASTROFF:} C = test real code with \texttt{unittest}; A = evaluator/maintainer; S = grouped tests + discovery; T = precise/repeatable; R = test-design assistant; O = runnable tests + evidence; F = Gamma behavior/errors; Function = satisfy P8 with automated tests.\\
\textbf{Output summary:} grouped test suite covering parsing, Gamma values, domain/range and user-facing service.\\
\textbf{Decision:} \textcolor{Green}{\bfseries Accepted} grouped suite; \textcolor{Teal}{\bfseries Revised} with explicit oracles and a clear 22-unit-test vs 8-validation-case distinction; \textcolor{Orange}{\bfseries Rejected} any use of \texttt{math.gamma()} as production logic.
\end{alertblock}

\vspace{0.14cm}
\begin{block}{D3 requirement coverage}
\scriptsize
\textbf{P7:} PEP 8 + Flake8 snapshot; pdb snapshot + conclusion; Pylint snapshot + remaining-message rationale; Semantic Versioning; radial UIDP mind map using the course principles; Accessibility Features and Evidence.\\
\textbf{P8:} 22 PyUnit tests + framework execution; representative cases/oracles; separate eight-case numerical validation.\\
\textbf{Consistency check:} final repository snapshot was re-run locally: \textbf{22/22 unit tests passed} and the \textbf{eight-case validation passed with max error $1.332\times10^{-15}$}.
\end{block}

\vspace{0.14cm}
\begin{block}{References}
\small
[1] SOEN 6011, \emph{Project Description}, Summer 2026; [2] P. Kamthan, \emph{User Interface Design Principles}, course notes, 2026; [3] Python Software Foundation, \emph{PEP 8}; [4] PyCQA, \emph{Flake8}; [5] PyCQA, \emph{Pylint}; [6] Python Software Foundation, \emph{pdb}; [7] Python Software Foundation, \emph{unittest}.\\[0.05cm]
[8] \emph{Semantic Versioning 2.0.0}; [9] C. Lanczos, ``A precision approximation of the gamma function,'' 1964; [10] NIST DLMF, Ch. 5, \emph{Gamma Function} (including reflection/domain behavior); [11] Python \texttt{math.gamma} validation oracle; [12] H. Tavakoli, \emph{Prompt Engineering for Everyone}, Ch. 6: CASTROFF, 2026; [13] Dreuw and Deselaers, \emph{beamerposter}.
\end{block}
\end{column}
\end{columns}

\end{frame}
\end{document}
